\subsection{Actual hermitian isometry over global function fields}
\label{C:global-function-hermitian-subsection}

In this subsection $k$ is a global function field of odd
characteristic, supplied as a finite field tower over
$\mathbf F_q(t)$.  Finite extensions, multiplication tables,
and involutions are supplied by exact finite data.
We use the field arithmetic and enumeration of all places in
Lemma~\ref{C:global-function-field-arithmetic}.
In particular a completion $k_v$ has an explicitly presented
finite residue field, a uniformizer, and computable finite
Laurent expansions of every specified global element.
All searches below are ordinary dovetailed searches with
proved termination.  We impose no complexity bound.

\begin{lemma}[Proved: local quadratic and quaternion subform decisions]
\label{C:global-function-local-form-algorithms}
Let $F=k_v$ be any specified completion.
Isometry of nonsingular quadratic forms given over $k$,
and representation of a specified invertible pure quaternion
$\lambda$ by a nonsingular skew-hermitian form over a
quaternion algebra $(a,b)_k$ with canonical involution,
are decidable after extension to $F$.
The same quadratic procedures apply to the finite algebraic
local names constructed in the proof.

Over $k$ itself, quadratic isometry, isotropy, square
membership and norms from any specified quadratic field
extension are decidable.  A positive decision permits
construction of an actual global solution.
\end{lemma}
\begin{proof}
\emph{The local quadratic procedure.}
For a matrix over the exact global subfield, diagonalize
by ordinary congruence.  The auxiliary algebraic local
names use the certified finite-precision procedure below.
If necessary choose a vector of nonzero norm
from the standard basis vectors and their pairwise sums;
characteristic different from two guarantees that one
works.  For each nonzero diagonal coefficient $c$,
compute its valuation and divide by a square power of
the uniformizer.  The form becomes
\[
                       q_0\perp\langle\varpi\rangle q_1,
                                                               \tag{C.3.GH1}
\]
where both $q_0,q_1$ have unit diagonal coefficients.
Their reductions are explicit nonsingular forms over
the finite residue field.

We recall the exact valuation argument behind this
procedure, also used in
\cite[Section 2, p.\ 2]{C-DarkeyMensahGlobalForms2021}.
A nonzero isotropic residue vector of a unit form
lifts to an isotropic vector: at least one partial
derivative is a unit, so the simple-root form of
Hensel's lemma applies with the other coordinates fixed.
The resulting primitive vector splits a unimodular
hyperbolic plane.  Continuing gives the lifted residue
Witt decomposition.  If both remaining residue forms
are anisotropic, (C.3.GH1) is anisotropic.
Indeed, after normalizing a putative zero vector to
have a unit coordinate, the least valuation among
its even and odd blocks has one definite parity;
reduction of the terms of that valuation would give
a nonzero isotropic residue vector in that block.
This is impossible.
Thus (C.3.GH1) is hyperbolic precisely when both
residue forms are hyperbolic.
Over a finite field of odd order a nonsingular form
of dimension $2s$ is hyperbolic exactly when its
determinant is $(-1)^s$ times a square; odd-dimensional
forms are not hyperbolic.
This is a finite arithmetic test.
Two forms of equal dimension are isometric precisely
when their difference is hyperbolic, by Witt
cancellation.  These facts decide local isometry,
isotropy and the anisotropic dimension.
They also show that the four square classes of $F$
are represented by $1,u,\varpi,u\varpi$, where
$u$ lifts a residue nonsquare.

\emph{A local quaternion division algebra.}
Write $D=(a,b)_F$ with standard pure generators
$i,j$ satisfying $i^2=a$, $j^2=b$, $ij=-ji$.
Its norm form is
$\langle1,-a,-b,ab\rangle$.
The preceding isotropy test decides whether $D$
is split or is a division algebra.
In the division case the local isometry theorem
\cite[Theorem 2.4(i), pp.\ 6--7]{C-LewisUngerVanGeel2004}
classifies skew-hermitian forms by rank and
\[
        \operatorname{disc}(H)
          =(-1)^n\operatorname{Nrd}(H)\quad\bmod F^{\times2}
                                                               \tag{C.3.GH2}
\]
for an $n$ by $n$ Gram matrix.
All these data are computable.
For example a finite splitting field computes the
reduced norm by the determinant of the split matrix,
as in Lemma~\ref{C:global-function-central-simple-algorithms}.

There is a finite test for the possible discriminants
of a rank-one skew form; this avoids any convention
about a tabulated rank-one sign.
The square of $xi+yj+zij$ is
$ax^2+by^2-abz^2$.
For each of the four nonzero square classes $c$,
test isotropy of
\[
                         \langle a,b,-ab,-c\rangle .
                                                               \tag{C.3.GH3}
\]
The pure ternary form is anisotropic in a division
algebra, so an isotropic vector of (C.3.GH3) must
have nonzero last coordinate.
It therefore proves, and is equivalent to,
representation of $c$ by a pure square.
The resulting finite set is exactly the set of
rank-one discriminants in (C.3.GH2).
Skew-hermitian Gram--Schmidt diagonalization over
a quaternion division algebra shows that the
possible rank-$r$ discriminants are the products
of $r$ members of this finite set; for $r=0$
the only product is $1$.

Consequently a rank-$n$ form $h$ represents
$\lambda$ precisely when a rank-$(n-1)$
complement can have discriminant
$\operatorname{disc}(h)/\lambda^2$.
Existence of that complement is the finite
product test just described, and equality of
rank and discriminant proves its sufficiency.
A vector of invertible norm spans a nonsingular
free line, so this subform criterion is exactly
the requested representation criterion.

\emph{A split local quaternion algebra.}
We construct a matrix realization with computable
local coefficients.
If $a$ is a square in $F$, choose its Hensel square
root $r$ and use
\[
 i=\begin{pmatrix}r&0\\0&-r\end{pmatrix},
 \qquad j=\begin{pmatrix}0&b\\1&0\end{pmatrix}.
\]
Otherwise splitting of $(a,b)_F$ says that
$b=s^2-at^2$ for some $s,t\in F$.
These elements can be obtained without a negative
search for a local point.
Enumerate approximations in the dense subfield
$k$, rescaling the variables and the polynomial
by powers of $\varpi$ to make them integral.
Test the univariate Hensel inequality
$v(f(x_0))>2v(f'(x_0))$ with one coordinate fixed.
The nonzero norm equation is smooth, so at an
actual solution some partial derivative is
nonzero; sufficiently close approximations and
suitable rescalings satisfy this test.
The search therefore terminates.
Hensel iteration gives an exact algebraic local
root and its approximations to any finite order.
Now use
\[
 i=\begin{pmatrix}0&a\\1&0\end{pmatrix},
 \qquad
 j=\begin{pmatrix}s&-at\\t&-s\end{pmatrix}.
\]
The four matrices $1,i,j,ij$ give the required
isomorphism $D\simeq M_2(F)$.

Put $J=\left(\begin{smallmatrix}0&1\\-1&0\end{smallmatrix}\right)$.
Quaternion conjugation becomes
$X\mapsto J^{-1}X^{\mathsf t}J$.
If $\widetilde H$ is the ordinary $2n$ by $2n$
matrix of a skew-hermitian Gram matrix $H$,
then
$G=(I_n\otimes J)\widetilde H$ is symmetric
and nonsingular.
The identities
\[
  \widetilde{\bar P^{\mathsf t}HP}
      =(I_n\otimes J)^{-1}\widetilde P^{\mathsf t}
        G\widetilde P
\]
show directly that hermitian congruence is
quadratic congruence.
Representing $\lambda$ is therefore equivalent
to containing the nonsingular binary quadratic
form $J\widetilde\lambda$ as a subform of $G$.
List the finitely many diagonal forms of rank
$2n-2$ with entries among
$1,u,\varpi,u\varpi$.
The desired complement exists exactly when
one of their sums with $J\widetilde\lambda$
is isometric to $G$.
The local quadratic test decides this.

Only finitely many local digits are needed for
these tests, even when the splitting matrices
use the algebraic local names above.
One can enumerate changes of basis with
coefficients in the dense global subfield and
successively improve all approximations until
the change of basis and every leading
principal Gram minor are certified nonzero.
Such a basis exists: simultaneous nonvanishing
is a nonempty open condition for a nonsingular
symmetric form over an infinite field.
A nonzero minor is detected at a finite
precision.  Gaussian diagonalization in
that basis gives the valuations and leading
residue coefficients of every pivot.
These are exactly the finite data used in
(C.3.GH1); no decision from arbitrary
coefficient streams is being assumed.

\emph{Global decisions.}
The Hasse--Minkowski principles for quadratic
isotropy and hyperbolicity hold over the stated
global fields
\cite[Introduction, pp.\ 1--2]{C-LewisUngerVanGeel2004}.
Dovetail the following two searches.
The positive search enumerates exact global
isotropic vectors or isometry matrices and
checks their defining identities.
The negative search enumerates every place,
computes its local form, and stops at the
first local failure of isotropy or isometry.
The local procedures terminate, and the
local--global principles say that exactly
one of the searches succeeds.
For isometry first reject unequal dimensions;
Witt cancellation identifies the equal-dimension
case with hyperbolicity of the difference.
The zero-dimensional cases are direct.
This proves termination without requiring
any cutoff for the places to inspect.

Square membership is the binary isotropy
test for $\langle1,-c\rangle$, with $c=0$
handled separately.
For a quadratic field $k(\sqrt a)/k$,
membership of $c\ne0$ in its norm group
is equivalent to isotropy of
$\langle1,-a,-c\rangle$.
Because $a$ is nonsquare, a nonzero
isotropic vector has nonzero last
coordinate, so it gives the required
norm solution.  The positive search
constructs one.
This completes all assertions.
\end{proof}

\begin{theorem}[Proved: global quaternion skew-hermitian isometry is decidable]
\label{C:global-function-skew-hermitian-isometry}
Let $D$ be a quaternion division algebra over $k$
with canonical involution.  Isometry of arbitrary
nonsingular skew-hermitian forms over $D$ is
decidable by an ordinary terminating algorithm.
A positive decision constructs an actual
isometry.  The algorithm retains the
rank-one global norm obstruction, so it
does not assert the false local--global
principle for isometry of these forms.
\end{theorem}
\begin{proof}
Diagonalize both forms over $D$.
The usual orthogonalization is effective:
enumerate a vector with nonzero self-pairing,
split its nonsingular line, and repeat.
Such a vector exists for a nonzero
skew-hermitian form over a quaternion
division algebra.  Otherwise polarization
would make every value central and fixed
by conjugation while also allowing
multiplication in one argument by every
quaternion; nondegeneracy contradicts
the existence of noncentral quaternions.
Equivalently, if a nonzero off-diagonal
value remains, vectors $e_i+e_ja$
with $a$ ranging over a quaternion
basis exhibit a nonzero self-pairing.
All nonzero quaternion coefficients
are invertible, so the orthogonal
complement is computed by linear
algebra over $D$.

For rank $n\geq2$, let $\lambda$
be the first diagonal entry of the
target form $g$.
The representation theorem of
Kneser--Springer, in the precise form
\cite[Theorem 2.7(ii), p.\ 7]{C-LewisUngerVanGeel2004},
states that a nonsingular rank-$n$
skew-hermitian form represents this
invertible pure quaternion globally
if and only if it represents it at
every place.
Dovetail enumeration of exact global
vectors $v$ satisfying $h(v,v)=\lambda$
against the local nonrepresentation
tests of the preceding lemma over
all places.
The theorem proves that one search
terminates.
A local failure proves that $h$ and
$g$ are not isometric.
Otherwise split the constructed
global line and compare its
orthogonal complement with the
remaining target diagonal form.
Witt cancellation proves that this
comparison has exactly the same
answer as the original isometry
problem, independently of the
particular represented vector chosen.
Positive ranks decrease.

It remains to decide rank one;
write the two pure coefficients
as $\alpha,\beta\in D^\times$.
An isometry is an element $x\in D^\times$
with $\bar x\alpha x=\beta$.
Taking reduced norms shows that
\[
                    c^2=\operatorname{Nrd}(\beta)/
                              \operatorname{Nrd}(\alpha),
                    \qquad c=\operatorname{Nrd}(x).
                                                               \tag{C.3.GH4}
\]
The preceding lemma decides whether
the displayed ratio is a square
and constructs its two possible
square roots.  If it is nonsquare,
there is no isometry.
For each root $c$ put $\delta=\beta/c$.
Then $\delta^2=\alpha^2$.
Choose a nonzero $u$ satisfying
$\alpha u=u\delta$ as follows.
If $\alpha+\delta\ne0$, take
$u=\alpha+\delta$.
Otherwise choose a basis element
$b$ not commuting with $\alpha$
and put $u=\alpha b-b\alpha$.
It is nonzero and anticommutes
with $\alpha$, so it works when
$\delta=-\alpha$.
Such a $b$ exists because a
nonzero pure quaternion is
noncentral.
Every nonzero $u$ is invertible.

The centralizer of $\alpha$ is
the quadratic field $L=k(\alpha)$.
Its involution is the restriction
of quaternion conjugation.
For $z\in L^\times$ we have
\[
 \overline{zu}\,\alpha(zu)
       =N_{L/k}(z)\,
          \frac{\operatorname{Nrd}(u)}{c}\,\beta .
                                                               \tag{C.3.GH5}
\]
Consequently this root $c$
permits an isometry precisely when
$c/\operatorname{Nrd}(u)$ is a
norm from $L$.
For necessity, any isometry $x$
with reduced norm $c$ satisfies
$\alpha x=x\delta$; hence
$xu^{-1}\in L^\times$ and its
norm is the required ratio.
For sufficiency (C.3.GH5) gives
the isometry $x=zu$.
Both norm tests terminate by
the preceding lemma.
Testing both signs of $c$ is
essential.
This proves the rank-one decision
and completes the induction.
Composing the recorded basis
changes constructs a positive
isometry.
\end{proof}

\begin{theorem}[Proved: all residual hermitian arithmetic over global function fields]
\label{C:global-function-hermitian-isometry}
Let $D$ be a finite-dimensional division $k$-algebra,
with center a specified finite extension $K/k$,
and let $\sigma$ be a specified $k$-linear involution.
For either sign $\varepsilon$, nonsingular
$\varepsilon$-hermitian forms over $(D,\sigma)$
have decidable actual isometry and a literal
decidable least-code index set.
The additive Grothendieck group of their
orthogonal-sum monoid has decidable canonical
finite coordinates, computable addition and
inverse, and computable actual representatives
of every finite virtual input.
No division-algebra hermitian isometry
predicate remains as an oracle.
\end{theorem}
\begin{proof}
The center, its finite tower, and reduced norms
are computed by
Lemma~\ref{C:global-function-central-simple-algorithms}.
All needed linear systems are finite systems
over the supplied global field.

\emph{First kind.}
If $\sigma$ fixes $K$, it identifies $D$
with its opposite as a central $K$-algebra.
Thus $2[D]=0$ in the Brauer group.
Period equals index over every global field
\cite[Theorem 1.5.36(iv), pp.\ 25--26]
 {C-PoonenRationalPoints2017}.
Since $D$ is a division algebra, its degree
is therefore one or two.
In degree one, symmetric forms are the
quadratic case of the first lemma, and
alternating forms have their usual unique
isometry class in each even dimension.

In degree two compute quaternion conjugation
$\bar x=\operatorname{Trd}(x)-x$.
A nonzero pure element $i$ and a nonzero
commutator $j=ib-bi$ give
$i^2=a$, $j^2=b'$ in $K^\times$,
$ij=-ji$, and the quaternion basis
$1,i,j,ij$.
One finds them by finite linear algebra
and verifies the basis.
The automorphism $\sigma\circ\bar{\phantom{x}}$
is inner by Skolem--Noether
\cite[Lemma 11.6.2, Tag 074R]
 {C-StacksGlobalFormSkolem}.
Solve the linear equations
$\sigma(x)u=u\bar x$ on a basis
to find a nonzero implementer $u$.
It is invertible, and $\sigma^2=1$
implies $\bar u=\eta u$ for
$\eta\in\{1,-1\}$.
To check the last assertion explicitly,
inner conjugation by $u\bar u^{-1}$
is the identity, so that element
is central; applying conjugation
again makes its scalar square one.
The map
\[
                 H\longmapsto u^{-1}H
                                                               \tag{C.3.GH6}
\]
identifies $\varepsilon$-hermitian
congruence for $\sigma$ with
$\varepsilon\eta$-hermitian
congruence for quaternion conjugation.
This follows by multiplying the
Gram congruence identity by $u^{-1}$.
Thus the negative canonical sign
is exactly the preceding theorem.

For the positive canonical sign,
Gram--Schmidt gives diagonal
coefficients in $K^\times$.
Every such scalar is a reduced
norm from $D$.
Indeed the quaternion norm form
is anisotropic, while adjoining
$\langle-c\rangle$ gives a
five-dimensional quadratic form.
Over each completion its even
and odd residue blocks in
(C.3.GH1) have total dimension
five; one has dimension at least
three and is isotropic over the
finite residue field.
The Hasse principle therefore
makes this five-dimensional
form isotropic over $K$.
The last coordinate cannot be
zero because the quaternion
norm is anisotropic.
After division by that coordinate
we obtain an element of norm $c$.
Exact enumeration constructs it.
Rescaling every diagonal vector
therefore makes all coefficients
one.  Rank is the complete
invariant in this case.

\emph{Second kind.}
Suppose the restriction of $\sigma$
to $K$ is nontrivial, and let
$K_0=K^\sigma$.
Compute this field as the kernel
of $\sigma-1$ on the finite
$k$-vector space $K$.
Its multiplication table and a
finite tower are explicit.
The field $K_0$ is again global,
and $K/K_0$ is quadratic.
Choose $c\in K^\times$ with
$\sigma(c)=-c$ by subtracting
the conjugate of a nonfixed
basis element.
Then $K=K_0(c)$ and $c^2\in K_0$.
Multiplying a skew-hermitian
Gram matrix by $c$ turns it
into a hermitian one, preserving
all isometries.
For hermitian matrices $H,H'$
of equal rank, the global
classification is
\[
 H\simeq H'
 \quad\Longleftrightarrow\quad
 \frac{\operatorname{Nrd}(H)}
      {\operatorname{Nrd}(H')}
        \in N_{K/K_0}(K^\times).
                                                               \tag{C.3.GH7}
\]
This is exactly
\cite[Section 3, Theorem 3.1 and
Corollary 3.2, pp.\ 18--19]
 {C-LewisUngerVanGeel2004}:
its global-field hypothesis
includes the present function
fields, and there are no real
places requiring signature data.
The reduced norms in (C.3.GH7)
belong to $K_0^\times$.
They are computed determinants
in a finite splitting field,
so the norm test of the first
lemma over $K_0$ decides
(C.3.GH7).
If an explicit isometry is
required, enumerate matrices
over $D$ and check the exact
Gram and inverse identities;
the positive decision guarantees
termination.

\emph{Canonical finite coordinates.}
Code finite matrices using the
decidable field and algebra
names.  The hermitian equations
and invertibility are decidable,
and every isometry comparison
has just been made decidable.
Retain a valid matrix code
exactly when no smaller valid
code of the same rank and sign
is isometric to it.
There are finitely many
predecessors, giving a decidable
literal index set.
Every form has one least code,
computed by the same finite
predecessor tests.

Witt cancellation over division
algebras with involution and
$2\ne0$ makes the orthogonal-sum
monoid cancellative.
The hypotheses are the
finite-dimensional hermitian
ones in the already checked
transfer theorem
\cite[Section 3, pp.\ 271--273]{C-QSS1979}.
Represent a virtual class by
a pair $(H_+,H_-)$.
Two pairs have the same class
precisely when
\[
                 H_+\perp H'_-\simeq H'_+\perp H_- .
                                                               \tag{C.3.GH8}
\]
This is now a proved terminating
test, not an uncomputed quotient.
Retaining the least valid pair
code among its finitely many
predecessors gives canonical
finite group coordinates.
Orthogonal sum followed by
canonicalization computes addition;
interchanging the two matrices
computes inverse.
The chosen pair is itself an
actual positive and negative
representative.
These conclusions include both
signs and all division degrees.
\end{proof}
